% GENERATED, edit the French unit or the translation table
\mxunit{automatismes-ch08}{Variations and extrema of a function}

\begin{fluencysheet}{Thu~11~Feb.~2027}{Increasing, decreasing, monotonicity}
\mxfq{Warm-up · Quick}{IMG}{$f(x)=-2x+5$. Calculate $f(1)$ and $f(3)$; what do you notice?}{$f(1)=3$, $f(3)=-1$: the order of the images is the reverse of that of the inputs ($f$ is decreasing).}{}
\mxfq{Core · MCQ}{MOD}{$f$ is increasing on $I$ means that, for all $a<b$ in $I$:\mxchoices{$a\leqslant b$ only}{$f(a)\geqslant f(b)$}{$f(a)\leqslant f(b)$}{$f(a)=f(b)$}}{C.\ $f(a)\leqslant f(b)$}{Increasing: preserves the order; decreasing: reverses it.}
\mxfq{Core + · Graph}{REG}{Curve of $f(x)=x^2$ on $\intff{-2}{2}$:\par\smallskip\begin{tikzpicture}[scale=0.45,>=stealth]
\draw[very thin,gray!45] (-2,-1) grid (2,4);
\draw[->] (-2,0)--(2.3,0); \draw[->] (0,-1)--(0,4.3);
\foreach \i in {-2,-1,1,2}\draw (\i,.1)--(\i,-.1) node[below,font=\tiny]{\i};
\foreach \j in {-1,1,2,3,4}\draw (.1,\j)--(-.1,\j) node[left,font=\tiny]{\j};
\draw[thick,mxink,domain=-2:2,samples=60,smooth] plot (\x,{\x*\x});

\end{tikzpicture}\par Give the variations of $f$.}{Decreasing on $\intff{-2}{0}$, increasing on $\intff{0}{2}$.}{Expression used to draw the curve: $f(x)=x^2$.}
\mxfq{Review · Quick}{PRO}{Which multiplier corresponds to an increase of $\pc{15}$? To a decrease of $\pc{15}$?}{$\num{1,15}$; $\num{0,85}$.}{}
\mxfq{Thread · Quick}{DIV}{Is $457$ divisible by $3$? By $9$? Justify.}{Sum of digits $=16$, not a multiple of $3$: divisible by neither $3$ nor $9$.}{}
\end{fluencysheet}

\begin{fluencysheet}{Mon~1\textsuperscript{st}~Mar.~2027}{Table of variations; maximum and minimum on an interval}
\mxfq{Warm-up · Quick}{IMG}{Calculate $x^2$ for $x=-2$, $0$, $3$.}{$4$; $0$; $9$.}{}
\mxfq{Core · Method}{REG}{$f$ is decreasing on $\intff{-3}{1}$ with $f(-3)=5$, $f(1)=-1$, then increasing on $\intff{1}{4}$ with $f(4)=4$. Give the maximum and the minimum of $f$ on $\intff{-3}{4}$.}{Maximum $5$ reached at $x=-3$ ($5>4$); minimum $-1$ reached at $x=1$.}{Compare the values at the endpoints of the interval AND at the changes of variation.}
\mxfq{Core + · True/false}{REG}{True or false? Justify. The maximum of a function on $\intff{a}{b}$ is always reached at $a$ or at $b$.}{False. Counterexample: $f(x)=-x^2+4$ on $\intff{-2}{2}$: maximum $4$ reached at $0$.}{}
\mxfq{Review · Table}{PNP}{Is this table a proportionality table?\par\smallskip\mxarraystretch{1.15}\begin{tabular}{|c|c|c|c|}\hline
$x$ & 1 & 2 & 3 \\ \hline
$y$ & 5 & 9 & 13 \\ \hline
\end{tabular}\par }{No: $\dfrac51\neq\dfrac92$ (constant increase $+4$ but $y(0)\neq0$).}{}
\mxfq{Thread · Method}{GEO}{Is the triangle with sides $5$, $12$ and $13$ right-angled? Which property should be used?}{$5^2+12^2=169=13^2$: converse of Pythagoras' theorem, right-angled (right angle opposite the side $13$).}{}
\end{fluencysheet}

\begin{fluencysheet}{Tue~2~Mar.~2027}{Affine functions: rate of change and variations}
\mxfq{Warm-up · Quick}{IMG}{$f(x)=3x-1$. Calculate $\dfrac{f(4)-f(1)}{4-1}$.}{$\dfrac{11-2}{3}=3$: this is the slope $m$.}{}
\mxfq{Core · Quick}{MOD}{Direction of variation of $f(x)=-4x+7$?}{$m=-4<0$: $f$ is strictly decreasing on $\R$.}{}
\mxfq{Core + · Proof}{LIT}{Show that $f(x)=5x-2$ is increasing on $\R$.}{For $a<b$: $f(b)-f(a)=5(b-a)>0$, so $f(a)<f(b)$.}{Model write-up for the proof required by the syllabus.}
\mxfq{Review · Table}{LEC}{Ages of the members of a club.\par\smallskip\mxarraystretch{1.15}\begin{tabular}{|c|c|c|c|c|}\hline
Age & $10$ & $11$ & $12$ & $13$ \\ \hline
Frequency & 3 & 6 & 10 & 6 \\ \hline
\end{tabular}\par How many members are at most $11$ years old? What relative frequency is that?}{$3+6=9$ members; $\dfrac{9}{25}=\pc{36}$.}{}
\mxfq{Thread · Programme → formula}{REG}{Programme: multiply by $3$ then subtract $2$. Give $f(x)$ then $f(4)$.}{$f(x)=3x-2$; $f(4)=10$.}{}
\end{fluencysheet}

\begin{fluencysheet}{Thu~4~Mar.~2027}{Absolute value and square functions: sign and variations}
\mxfq{Warm-up · Quick}{CAL}{Calculate $|{-}5|$, $(-5)^2$ and $|2-7|$.}{$5$; $25$; $5$.}{}
\mxfq{Core · Method}{REG}{Give the table of variations of $x\mapsto|x|$ and its sign.}{Decreasing on $\intof{-\infty}{0}$, increasing on $\intfo{0}{+\infty}$; minimum $0$ at $0$; always non-negative.}{}
\mxfq{Core + · Reasoning}{CMP}{Compare $(-3)^2$ and $(-2)^2$ then $3^2$ and $2^2$. Is the square function increasing on $\R$?}{$9>4$ in both cases, but $-3<-2$: squaring reverses the order on $\intof{-\infty}{0}$; no, it is not increasing on $\R$.}{}
\mxfq{Review · Method}{PRB}{Two fair dice are rolled. Probability that the sum is $7$?}{$6$ favourable cases out of $36$: $\dfrac16$.}{}
\mxfq{Thread · Method}{ALG}{Programme: choose a number; add $1$; square it; subtract the square of the starting number. What is the result for $6$? What is the general expression?}{For $6$: $49-36=13$. In general $(x+1)^2-x^2=2x+1$.}{}
\end{fluencysheet}

\begin{fluencysheet}{Mon~8~Mar.~2027}{Reciprocal function: variations; relative position of $y=x$ and $y=x^2$}
\mxfq{Warm-up · Quick}{CAL}{Calculate $\dfrac12$, $\dfrac1{-4}$ and $\dfrac1{\num{0,1}}$.}{$\num{0,5}$; $-\num{0,25}$; $10$.}{}
\mxfq{Core · Method}{CMP}{Compare $\dfrac13$ and $\dfrac15$; what can be said about the variations of the reciprocal function on $\intoo{0}{+\infty}$?}{$\dfrac13>\dfrac15$ although $3<5$: the reciprocal function is decreasing on $\intoo{0}{+\infty}$ (and on $\intoo{-\infty}{0}$, but not on $\R^*$).}{}
\mxfq{Core + · True/false}{CMP}{True or false? Justify. For every real number $x>0$, $x^2\geqslant x$.}{False. Counterexample: $x=\dfrac12$, $x^2=\dfrac14<\dfrac12$. In fact $x^2\geqslant x$ if and only if $x\geqslant1$ (for $x>0$).}{}
\mxfq{Review · Quick}{PER}{What are the front view and the top view of a cylinder standing on its base?}{Front: a rectangle; top: a disc.}{}
\mxfq{Thread · Prove}{LIT}{Show that the product of two even integers is a multiple of $4$.}{With $a=2p$, $b=2q$: $ab=4pq$, a multiple of $4$.}{}
\end{fluencysheet}

\begin{fluencysheet}{Tue~9~Mar.~2027}{Comparing $f(a)$ and $f(b)$; optimisation problems}
\mxfq{Warm-up · Quick}{CMP}{$f(x)=x^2$. Compare $f(-3)$ and $f(2)$.}{$f(-3)=9>4=f(2)$.}{}
\mxfq{Core · Method}{REG}{$f(x)=x^2$, $a=-\num{1,5}$ and $b=-2$. Compare $f(a)$ and $f(b)$ without calculating, using the variations.}{$-2<-\num{1,5}\leqslant0$ and $f$ is decreasing on $\intof{-\infty}{0}$: $f(-2)>f(-\num{1,5})$. Check: $4>\num{2,25}$.}{}
\mxfq{Core + · Optimisation}{MOD}{$A(x)=x(10-x)$ for $x\in\intff{0}{10}$ (area of a rectangle with perimeter $20$). Show that $A(x)=25-(x-5)^2$ and deduce the maximum.}{$25-(x-5)^2=25-x^2+10x-25=x(10-x)$; maximum $25$ reached for $x=5$ (square of side $5$).}{}
\mxfq{Review · Quick}{STA}{$18$ pupils out of $24$ have an average above $10$. Give the relative frequency as a percentage.}{$\dfrac{18}{24}=\pc{75}$.}{}
\mxfq{Thread · MCQ}{EQ}{The solution of $\dfrac{x}{3}+1=\dfrac{7}{2}$ is:\mxchoices{$\dfrac{15}{2}$}{$\dfrac{9}{2}$}{$\dfrac{21}{2}$}{$\dfrac{5}{6}$}}{A.\ $\dfrac{15}{2}$}{Pitfall: forgetting to subtract the $1$ before multiplying by $3$ ($\frac{21}{2}$).}
\end{fluencysheet}
